\documentclass[11pt,a4paper,oneside]{report}
\usepackage[spanish,activeacute,es-noshorthands]{babel}
\usepackage[utf8]{inputenc}
\usepackage[T1]{fontenc}
\usepackage{lmodern}
\usepackage{graphicx}
\usepackage[dvipsnames]{xcolor}
\usepackage{geometry}
\geometry{margin=2.3cm, headheight=14pt}
\usepackage{microtype}
\usepackage{booktabs}
\usepackage{longtable}
\usepackage{enumitem}
\usepackage{fancyhdr}
\usepackage{tcolorbox}
\tcbuselibrary{skins, breakable}
\usepackage{caption}
\captionsetup{font=small, labelfont=bf}
\usepackage{titlesec}
\usepackage{hyperref}
\hypersetup{
  colorlinks=true,
  linkcolor=NavyBlue,
  citecolor=ForestGreen,
  urlcolor=BrickRed,
  pdftitle={Bibliografia --- IA en operaciones espaciales},
  pdfauthor={CubeSat-EdgeAI-Payload}
}

\definecolor{intisat}{RGB}{26, 54, 93}
\definecolor{intisataccent}{RGB}{197, 143, 0}

\titleformat{\chapter}[hang]
  {\normalfont\Huge\bfseries\color{intisat}}
  {\thechapter.}{0.6em}{}
\titlespacing*{\chapter}{0pt}{-20pt}{20pt}

\titleformat{\section}[hang]
  {\normalfont\Large\bfseries\color{intisat}}
  {\thesection}{0.6em}{}

\pagestyle{fancy}
\fancyhf{}
\fancyhead[L]{\nouppercase{\leftmark}}
\fancyhead[R]{\thepage}
\renewcommand{\headrulewidth}{0.4pt}

\newtcolorbox{nicho}[1][]{
  colback=yellow!7, colframe=intisataccent,
  fonttitle=\bfseries, title=Nicho del proyecto CubeSat-EdgeAI-Payload,
  breakable, sharp corners=south, #1
}

\begin{document}

% ─── Portada ────────────────────────────────────────────────────────
\begin{titlepage}
\centering
\vspace*{1cm}
{\color{intisat}\rule{\textwidth}{2pt}}
\vspace{0.6cm}

{\Huge\bfseries\color{intisat} Inteligencia Artificial\\[0.2em]
en operaciones espaciales}

\vspace{0.4cm}
{\Large\bfseries Bibliografia comentada y panorama actual}

\vspace{0.7cm}
{\Large\itshape Misiones, hardware, papers, empresas y datasets\\
para procesamiento autonomo a bordo de satelites y rovers}

\vspace{1cm}
{\color{intisataccent}\rule{0.6\textwidth}{1pt}}

\vspace{2cm}

\begin{tabular}{rl}
\textbf{Proyecto:} & CubeSat-EdgeAI-Payload \\
\textbf{Plataforma:} & INTISAT --- payload de microscopia FPM \\
\textbf{Repositorio:} & \texttt{github.com/JaminYC/CubeSat-EdgeAI-Payload} \\
\textbf{Fecha:} & \today \\
\textbf{Documento:} & v1.0 --- compilacion inicial
\end{tabular}

\vfill

\begin{minipage}{0.85\textwidth}
\textbf{Resumen.}\quad
Compilacion de misiones reales, hardware certificado, empresas activas,
papers academicos, frameworks open-source, datasets publicos y conferencias
relevantes para el desarrollo de inteligencia artificial a bordo de
satelites. Cubre desde Earth Observation (PhiSat) hasta exploracion
planetaria (rovers de Mars), pasando por plataformas reconfigurables
(OPS-SAT) y robotica espacial (Astrobee, CIMON). Cada referencia incluye
la mision, el año, la aplicacion concreta y el hardware utilizado, con
el objetivo de servir como mapa de orientacion para proyectos de IA
embarcada en CubeSats academicos y comerciales.
\end{minipage}

\vspace{1.2cm}
{\color{intisat}\rule{\textwidth}{2pt}}
\end{titlepage}

\tableofcontents
\clearpage

% =====================================================================
\chapter{Observacion de la Tierra (EO)}

El area mas activa de IA on-board: filtros de nubes, deteccion de cambios,
clasificacion de cultivos, monitoreo de incendios. La mayoria de las
misiones que actualmente vuelan IA en orbita pertenecen a esta categoria.

\begin{longtable}{p{3.2cm}p{1cm}p{4.5cm}p{3.2cm}p{2.5cm}}
\toprule
\textbf{Mision} & \textbf{Año} & \textbf{Que IA hace} & \textbf{Hardware} & \textbf{Link} \\
\midrule
\endhead
PhiSat-1 (ESA)         & 2020      & filtra nubes con CNN antes de bajar imagenes & Intel Movidius Myriad 2 & \href{https://www.esa.int/Applications/Observing_the_Earth/Ph-sat}{esa.int/Ph-sat} \\
PhiSat-2 (ESA)         & 2024      & 6 apps simultaneas (cambio detectado, water quality) & Intel Myriad X & \href{https://www.esa.int/Applications/Observing_the_Earth/Phsat-2}{esa.int/Phsat-2} \\
HYPSO-1, HYPSO-2 (NTNU Noruega) & 2022, 2024 & ML hyperspectral para floraciones de algas & RPi CM4 & \href{https://www.ntnu.edu/ie/smallsat}{ntnu.edu} \\
Intuition-1 (KP Labs)  & 2023      & IA hyperspectral on-board & Antelope FPGA & \href{https://kplabs.space/intuition-1/}{kplabs.space} \\
CogniSAT-6 (Aiko + Open Cosmos) & 2024 & primer satelite EO comercial AI integrada & Movidius + custom & \href{https://aikospace.com/cognisat-6/}{aikospace.com} \\
Pleiades Neo (Airbus)  & 2021+     & clasificacion + super-resolucion & propietario & \href{https://www.intelligence-airbusds.com/pleiades-neo}{Airbus DS} \\
WorldView Legion (Maxar) & 2024    & deteccion de objetos a bordo & propietario & \href{https://www.maxar.com}{maxar.com} \\
\bottomrule
\end{longtable}

% =====================================================================
\chapter{Plataformas reconfigurables (AI playground en orbita)}

Satelites pensados para que cualquiera suba experimentos de IA --- el
equivalente a un \emph{lab de pruebas} en orbita.

\begin{longtable}{p{3.5cm}p{2.2cm}p{6cm}p{3.5cm}}
\toprule
\textbf{Mision} & \textbf{Periodo} & \textbf{Para que sirve} & \textbf{Link} \\
\midrule
\endhead
OPS-SAT (ESA)              & 2019--2024 & satelite plataforma para experimentos AI & \href{https://www.esa.int/Enabling_Support/Operations/OPS-SAT}{esa.int/OPS-SAT} \\
OPS-SAT VOLT (ESA)         & 2024+      & sucesor de OPS-SAT  & \href{https://www.esa.int/Enabling_Support/Operations/OPS-SAT_VOLT}{esa.int/VOLT} \\
D-Orbit ION SCV            & 2020+      & "satelite hostel" con AI a bordo & \href{https://www.dorbit.space/ion-scv}{dorbit.space} \\
NASA STARLING              & 2023       & swarm autonomo de 4 CubeSats con AI distribuida & \href{https://www.nasa.gov/smallsat-institute/sst-soa/starling}{nasa.gov} \\
\bottomrule
\end{longtable}

% =====================================================================
\chapter{Operaciones autonomas del satelite}

\begin{longtable}{p{3.2cm}p{4cm}p{6cm}p{2.5cm}}
\toprule
\textbf{Aplicacion} & \textbf{Mision / proyecto} & \textbf{Que hace} & \textbf{Link} \\
\midrule
\endhead
Anomaly detection  & OPS-SAT, NASA cFS-AI         & detecta fallos del bus antes que el OBC tradicional & \href{https://www.esa.int/Enabling_Support/Operations/OPS-SAT}{ESA} \\
Mission planning   & Aiko Mission Control AI      & distribuye tareas en constelacion automaticamente & \href{https://aikospace.com}{aikospace.com} \\
Self-healing FDIR  & NASA Resilient Spacecraft    & reconfiguracion autonoma ante fallos & \href{https://www.nasa.gov}{nasa.gov} \\
Cognitive radio    & NASA SCaN Testbed (ISS)      & aprende interferencia y ajusta frecuencias & \href{https://www.nasa.gov/scan}{nasa.gov/scan} \\
\bottomrule
\end{longtable}

% =====================================================================
\chapter{Mars y exploracion planetaria}

\begin{longtable}{p{3.5cm}p{4cm}p{8cm}}
\toprule
\textbf{Mision} & \textbf{IA on-board} & \textbf{Detalle} \\
\midrule
\endhead
Curiosity (NASA, 2012-) & AEGIS desde 2016 & la rover elige autonomamente que rocas estudiar con ChemCam \\
Perseverance (NASA, 2021-) & AutoNav & conduccion autonoma con vision a 5 m/s. Triplica el ritmo vs teleop \\
Ingenuity helicopter (2021-2024) & navegacion vision-only & SLAM completo a bordo con IMU + camara \\
PIXL (Perseverance) & adaptive raster scanning con ML & optimiza secuencia de medidas espectrales \\
\bottomrule
\end{longtable}

\textbf{Links}:
\begin{itemize}[leftmargin=2em]
  \item AEGIS paper: \url{https://www.science.org/doi/10.1126/scirobotics.aan4582}
  \item AutoNav: \url{https://mars.nasa.gov/mars2020/mission/technology/}
  \item Ingenuity: \url{https://www.jpl.nasa.gov/news/ingenuity-tech-demo}
\end{itemize}

% =====================================================================
\chapter{ISS — robots con IA}

\begin{longtable}{p{3cm}p{3cm}p{6cm}p{3cm}}
\toprule
\textbf{Robot} & \textbf{Quien} & \textbf{Que hace} & \textbf{Link} \\
\midrule
\endhead
Astrobee (NASA)        & desde 2019 & 3 robots free-flying con vision + planeamiento autonomo & \href{https://www.nasa.gov/astrobee}{nasa.gov/astrobee} \\
CIMON-2 (DLR Alemania) & desde 2019 & asistente AI de voz para astronautas, basado en Watson & \href{https://www.airbus.com/en/cimon}{airbus.com} \\
Robonaut 2 (NASA)      & 2011--2018 & manipulacion humanoide con vision profunda & \href{https://robonaut.jsc.nasa.gov}{nasa.gov} \\
Astro Pi (ESA + RPi Foundation) & desde 2015 & dos Pi 4 en la ISS para experimentos estudiantiles + IA & \href{https://astro-pi.org}{astro-pi.org} \\
\bottomrule
\end{longtable}

% =====================================================================
\chapter{Rendezvous, docking, landing autonomo}

\begin{longtable}{p{4cm}p{2.5cm}p{8cm}}
\toprule
\textbf{Mision} & \textbf{Año} & \textbf{Aplicacion} \\
\midrule
\endhead
Dragon (SpaceX)              & 2012-      & docking automatico con vision + LIDAR \\
Soyuz KURS-NA                & desde 2019 & estimacion de relativa con ML \\
OSIRIS-REx (NASA)            & 2020       & TAGSAM: maniobra autonoma para tocar el asteroide Bennu \\
HERA (ESA)                   & 2024-2026  & navegacion autonoma alrededor de Didymos \\
Mars Sample Return (NASA)    & 2030+      & rendezvous + docking en orbita marciana, todo autonomo \\
\bottomrule
\end{longtable}

% =====================================================================
\chapter{Astronomia y ciencia espacial}

\begin{longtable}{p{4cm}p{10cm}}
\toprule
\textbf{Telescopio} & \textbf{Que IA usa} \\
\midrule
\endhead
JWST (NASA)                          & pipeline ML para clasificacion estelar y artefacto removal \\
Euclid (ESA)                         & clasificacion de morfologias galacticas con ML \\
Kepler / TESS                        & deteccion de exoplanetas con CNN (post-procesado) \\
Roman Space Telescope (NASA, 2027)   & traccion de microlensing con IA tiempo real \\
\bottomrule
\end{longtable}

% =====================================================================
\chapter{Hardware AI usado en orbita}

\begin{longtable}{p{4.5cm}p{4cm}p{6cm}}
\toprule
\textbf{Chip} & \textbf{Donde se usa} & \textbf{Notas} \\
\midrule
\endhead
Intel Movidius Myriad 2 / X    & PhiSat, CogniSAT       & el caballito de batalla. $\sim 1$ W. 1 TOPS \\
Google Edge TPU (Coral)        & experimentos en LEO    & testeado para radiacion por NASA y JAXA \\
NVIDIA Jetson Nano/Orin        & algunos commercial sats & mas potente pero 5-10 W \\
AMD Zynq Ultrascale+ MPSoC     & Mars rovers, JWST      & radiation-hardened version \\
Brainchip Akida (neuromorphic) & testeo en orbita 2024  & spike-based, ultra bajo consumo \\
KP Labs Antelope/Leopard       & Intuition-1            & FPGA + AI accelerators custom \\
Ubotica CogniSAT-XE2           & varios                 & computer vision dedicado \\
RPi 4 / CM4                    & HYPSO-1, Astro Pi      & COTS, no rad-hard, vida 1-3 años \\
\bottomrule
\end{longtable}

% =====================================================================
\chapter{Empresas y actores principales}

\begin{longtable}{p{4cm}p{6cm}p{4cm}}
\toprule
\textbf{Empresa} & \textbf{Que hace} & \textbf{Link} \\
\midrule
\endhead
ESA $\Phi$-lab            & division de IA de ESA & \href{https://philab.esa.int}{philab.esa.int} \\
NASA Frontier Development Lab & sprints anuales de IA aplicada al espacio & \href{https://frontierdevelopmentlab.org}{frontierdevelopmentlab.org} \\
Ubotica (Irlanda)         & hardware AI para PhiSat & \href{https://ubotica.com}{ubotica.com} \\
KP Labs (Polonia)         & Antelope edge AI processor & \href{https://kplabs.space}{kplabs.space} \\
Aiko Space (Italia)       & software AI mision control & \href{https://aikospace.com}{aikospace.com} \\
Cognitive Space (USA)     & constelaciones autonomas & \href{https://cognitivespace.com}{cognitivespace.com} \\
Open Cosmos (UK)          & satelite plataforma + AI & \href{https://www.open-cosmos.com}{open-cosmos.com} \\
Spaceability AI (UK)      & dashboards AI para operadores & \href{https://spaceability.com}{spaceability.com} \\
LeoLabs (USA)             & SDA con ML & \href{https://leolabs.space}{leolabs.space} \\
Slingshot Aerospace (USA) & digital twin de satelites con AI & \href{https://slingshot.space}{slingshot.space} \\
D-Orbit (Italia)          & ION platform & \href{https://www.dorbit.space}{dorbit.space} \\
Loft Orbital (USA)        & mision como servicio & \href{https://loftorbital.com}{loftorbital.com} \\
Spire Global (USA)        & constelacion 100+ CubeSats & \href{https://spire.com}{spire.com} \\
Planet Labs (USA)         & Doves, 200+ CubeSats EO & \href{https://www.planet.com}{planet.com} \\
\bottomrule
\end{longtable}

% =====================================================================
\chapter{Frameworks y herramientas open-source}

\begin{longtable}{p{4cm}p{6cm}p{4cm}}
\toprule
\textbf{Tool} & \textbf{Para que} & \textbf{Link} \\
\midrule
\endhead
NASA cFS                  & sistema operativo de vuelo + extensiones AI & \href{https://github.com/nasa/cFS}{github/nasa/cFS} \\
NASA F'                   & framework usado en MarCO, Ingenuity & \href{https://github.com/nasa/fprime}{github/nasa/fprime} \\
ESA Onboard AI demos      & repositorio publico & \href{https://gitlab.esa.int/PhiLab}{gitlab.esa.int} \\
OpenVINO                  & wrapper para deploy en Movidius & \href{https://github.com/openvinotoolkit/openvino}{openvinotoolkit} \\
TensorFlow Lite Micro     & inferencia en MCU constrained & \href{https://github.com/tensorflow/tflite-micro}{tflite-micro} \\
OreSat (PSAS)             & CubeSat 100\% open: hardware, firmware, ground station & \href{https://github.com/oresat}{github/oresat} \\
UPSat (LSF Greece)        & otro CubeSat 100\% open con docs operativas & \href{https://gitlab.com/librespacefoundation/upsat}{gitlab.com} \\
SatNOGS DBs               & datasets de telemetria & \href{https://db.satnogs.org}{db.satnogs.org} \\
\bottomrule
\end{longtable}

% =====================================================================
\chapter{Papers de referencia (reviews)}

\begin{longtable}{p{8cm}p{1.5cm}p{6cm}}
\toprule
\textbf{Titulo} & \textbf{Año} & \textbf{Por que vale} \\
\midrule
\endhead
\emph{AI in Space: Opportunities and Challenges} (Furano et al., ESA) & 2020 & reporte oficial de ESA con todo el panorama \\
\emph{A Survey of On-Board AI for Space Applications} (Salvo et al., IEEE) & 2023 & review tecnico exhaustivo \\
\emph{Onboard AI for Earth Observation} (Esposito et al.) & 2022 & foco en EO especificamente \\
\emph{Lessons from PhiSat-1} (Marcuccio et al.) & 2021 & post-mortem del primer satelite IA \\
\emph{Edge ML for Space Systems} (Kabir et al.) & 2024 & hardware comparison para edge ML en orbita \\
\emph{Cognitive Space Systems: state-of-the-art and challenges} (Ortiz et al.) & 2023 & autonomia, FDIR, cognitive radio \\
\emph{State of the Art of CubeSats} (Bouwmeester et al.) & 2010+ & el clasico panorama general \\
\emph{Why CubeSats Fail} (Swartwout, IEEE Aerospace) & 2018 & top 10 razones reales de fallos \\
\bottomrule
\end{longtable}

\textbf{Buscar texto completo en}:
\begin{itemize}[leftmargin=2em]
  \item \url{https://arxiv.org}
  \item \url{https://digitalcommons.usu.edu/smallsat/}
\end{itemize}

% =====================================================================
\chapter{Conferencias y workshops dedicados}

\begin{longtable}{p{4.5cm}p{6cm}p{4cm}}
\toprule
\textbf{Evento} & \textbf{Foco} & \textbf{Link} \\
\midrule
\endhead
OBDP (On-Board Data Processing) & bianual ESA, lo mas alineado a IA on-board & \href{https://obdp2024.sciencesconf.org}{obdp2024} \\
AI4Space (CVPR/ICCV)          & computer vision para espacio & \href{https://aiforspace.github.io}{aiforspace.github.io} \\
SmallSat USU                  & el grande, todo lo de SmallSat & \href{https://digitalcommons.usu.edu/smallsat}{digitalcommons.usu.edu} \\
Space Computing Conference (NASA) & hardware a bordo & \href{https://nepp.nasa.gov}{nepp.nasa.gov} \\
EDHPC                         & European Data Handling and Data Processing & \href{https://www.edhpc-conference.com}{edhpc-conference.com} \\
IAC                           & International Astronautical Congress & \href{https://www.iafastro.org}{iafastro.org} \\
IEEE Aerospace Conference     & papers tecnicos & \href{https://www.aeroconf.org}{aeroconf.org} \\
\bottomrule
\end{longtable}

% =====================================================================
\chapter{Datasets publicos para entrenar IA espacial}

\begin{longtable}{p{4cm}p{6cm}p{5cm}}
\toprule
\textbf{Dataset} & \textbf{Que tiene} & \textbf{Link} \\
\midrule
\endhead
Sentinel-2                  & imagenes EO multiespectral & \href{https://scihub.copernicus.eu}{scihub.copernicus.eu} \\
Landsat                     & imagenes EO 50 años & \href{https://earthexplorer.usgs.gov}{earthexplorer.usgs.gov} \\
PROBA-V Cloud Detection     & el dataset que entreno PhiSat-1 & \href{https://kelvins.esa.int/proba-v-super-resolution/}{kelvins.esa.int} \\
Mars Curiosity images       & NASA Planetary Data System & \href{https://pds.nasa.gov}{pds.nasa.gov} \\
SatNOGS Network             & telemetria CubeSat & \href{https://network.satnogs.org}{network.satnogs.org} \\
xView2 / xBD                & damage assessment EO & \href{https://xview2.org}{xview2.org} \\
EuroSAT                     & clasificacion de cobertura terrestre & \href{https://github.com/phelber/EuroSAT}{github/EuroSAT} \\
BigEarthNet                 & Sentinel-2 con labels & \href{https://bigearth.net}{bigearth.net} \\
\bottomrule
\end{longtable}

% =====================================================================
\chapter{Donde se ubica CubeSat-EdgeAI-Payload}

El proyecto entra en la categoria de \textbf{Earth Observation con IA
on-board}, pero con la diferencia significativa de que aplica
\emph{microscopia biologica} en lugar de imagenes terrestres, lo que lo
ubica en un nicho menos poblado y diferenciado.

\section{Antecedentes directos}

\begin{itemize}[leftmargin=2em]
  \item \textbf{PharmaSat / GeneSat} (NASA Ames, 2006--2009): microscopia
        biologica en CubeSat sin IA on-board. Capturan imagenes y bajan
        crudo. Es el antecedente metodologico del experimento biologico.
  \item \textbf{PhiSat-1} (ESA, 2020): IA en orbita por primera vez, pero
        para EO. Establece el paradigma de procesar antes de bajar.
  \item \textbf{HYPSO-1} (NTNU, 2022): RPi CM4 con IA on-board, modelo
        similar de hardware al planeado para este proyecto.
\end{itemize}

\section{Aporte diferencial}

Pipeline IA completo (segmentacion, mejora, medicion) en payload de
microscopia biologica, no solo captura. Combina la linea de NASA Ames
(biologia en CubeSat) con el paradigma de PhiSat (IA on-board EO).

\begin{nicho}
\textbf{Nicho diferencial frente a la competencia actual}:
\begin{itemize}[leftmargin=1.5em]
  \item Imaging \emph{biologico}, no observacion terrestre.
  \item Lensless con OLED programable (apertura sintetica), poco usado en mision.
  \item Pipeline IA con multiples modelos (StarDist + Real-ESRGAN + N2V) en un solo Pi.
  \item Open-source (a diferencia de los desarrollos comerciales).
\end{itemize}
\end{nicho}

\section{Lo que falta para alcanzar nivel SOTA}

\begin{enumerate}[leftmargin=2em]
  \item Migrar a hardware certificado para vuelo (Movidius o Antelope, no
        RPi 5 de desarrollo). Para vuelo se planea \emph{Compute Module 5}.
  \item Tests TVAC, vibracion y radiacion segun GEVS / ECSS.
  \item Demostrar enlace uplink/downlink real con ground station.
  \item Compararse con un sistema de referencia comercial (Ubotica
        CogniSAT-XE2 o similar).
\end{enumerate}

% =====================================================================
\chapter{Lectura priorizada}

Si el tiempo es limitado, leer en este orden:

\begin{enumerate}[leftmargin=2.2em, label=\arabic*.]
  \item \textbf{PhiSat-1 paper} (Marcuccio 2021) --- el caso emblematico
        mas cercano al proyecto.
  \item \textbf{Furano et al. AI in Space} (ESA review) --- el panorama
        completo en un PDF.
  \item \textbf{OBDP 2024 proceedings} --- lo mas actual.
  \item \textbf{HYPSO papers} (NTNU) --- usan RPi CM4 con AI, hw mas
        cercano al nuestro.
  \item \textbf{AI4Space workshop papers} --- repositorio de ideas + benchmarks.
  \item \textbf{Why CubeSats Fail} (Swartwout) --- para evitar errores tipicos.
  \item \textbf{Bouwmeester State of the art of CubeSats} --- fundamentos.
\end{enumerate}

% =====================================================================
\chapter*{Como mantener este documento}
\addcontentsline{toc}{chapter}{Como mantener este documento}

Cuando alguien del equipo descubra una referencia nueva relevante:

\begin{enumerate}[leftmargin=2em]
  \item Identificar a que capitulo pertenece (1--13).
  \item Agregar fila a la tabla con fecha, descripcion y link.
  \item Si es un paper especialmente importante, sumarlo al ranking del
        capitulo de lectura priorizada.
  \item Compilar el \texttt{.tex} y commit + push.
\end{enumerate}

Si la referencia es de un area nueva no contemplada, abrir un capitulo
nuevo manteniendo el formato consistente.

\end{document}
